
Analysis of language model errors on benchmark datasets conventionally requires re-running models on test items to obtain item-level or category-level outputs. This approach creates barriers: API costs for frontier models can reach thousands of dollars for comprehensive evaluations, and access to some models is restricted or unavailable. When API credentials are unavailable, researchers may resort to synthetic data, which produces unreliable results.

An alternative approach would leverage published benchmark results from model technical reports. If these reports contain category-level breakdowns with sufficient granularity, they could support error analysis without re-evaluation. However, no prior work has systematically validated the availability, granularity, and completeness of such data across multiple model families and benchmarks.

This study addresses the question: Do published technical reports from major LLM labs contain category-level error rate data with sufficient granularity and completeness for systematic analysis? We focus on three model families (GPT from OpenAI, Claude from Anthropic, Llama from Meta), two benchmarks (TruthfulQA for truthfulness evaluation, MMLU for knowledge assessment), and two timepoints per family (baseline and current models).

We establish quantitative criteria for data sufficiency: (1) coverage across at least 3 model families with data for both baseline and current generations, (2) at least 10 categories per benchmark to provide adequate granularity, and (3) at least 90\% completeness to ensure statistical robustness. These thresholds derive from requirements for weak supervision approaches in related literature.

Our validation shows that all three model families publish category-level data meeting or exceeding all criteria. TruthfulQA results include 12 categories, MMLU results include 15 categories, and data completeness is 100\% using curated manual extraction. This establishes that category-level data are systematically available across major labs.

The contribution is limited to data availability validation. We do not test whether this data supports specific downstream applications such as error taxonomy generation, pattern clustering, or predictive modeling. Those applications remain as potential future work.
